\subsection{Controller 1 with Integral APF}

To achieve the control objective, the following error variables and auxiliary variables are first defined.
The position error \( {q}_{i0} = q_i - q_0 \) represents the position deviation between the follower and the target.
The velocity error: \( \dot{\tilde{q}}_i = \dot{q}_i - \dot{q}_0 \), representing the velocity deviation between the follower and the target.

First, we define an auxiliary variable as:
\begin{equation}\label{eq:sliding_var}
	s_i = \dot{q}_i - \zeta_i,
\end{equation}
where $\zeta_i = \dot{q}_0 - k_{\alpha} {q}_{i0} + \int_0^t  -k_{\alpha} {q}_{i0} + \phi_i(\tau) d \tau$.

Then, the control input \( u_i \) for the \( i \)-th follower is designed as:
\begin{equation}\label{eq:control_input}
	\tau_i = -k s_i + Y_i(q_{i},\dot{q}_{i},\dot{\zeta}_i,\zeta_i) \hat{\Theta}_i 
\end{equation}
where \( k > 0 \) and \( Y_i \): regression matrix, used to separate unknown parameters from known signals, and its expression is given by the following formula:
	\begin{equation}\label{eq:regression}
		M_i(q_i)\dot{\zeta}_i + C_i(q_i,\dot{q})\zeta_i + g_i(q_i)=Y_{i}(q_{i},\dot{q}_{i},\dot{\zeta}_i,\zeta_i)\Theta_{i}.
	\end{equation}
The term \( \hat{\theta}_i \) is the estimate of the unknown parameter \( \Theta_i \)  of  each agent;
To achieve accurate estimation of the unknown parameter \( \Theta_i \), the following parameter adaptation law is designed:
\begin{equation}\label{eq:adaptation}
	\dot{\hat{\Theta}}_i = -\Lambda^{-1} Y_i^T(q_{i},\dot{q}_{i},\dot{\zeta}_i,\zeta_i) s_i, 
\end{equation}
where \( \Lambda > 0 \) is the adaptation gain matrix, used to adjust the convergence speed of parameter estimation.

\begin{theorem}
	Under Assumptions A1--A3 and A4, the HELS \eqref{eq hetero agent} governed by \eqref{eq:control_input} and \eqref{eq:regression} fences the target in the sense that:
	\begin{enumerate}
		\item[\textbf{P1}] $\limsup_{t\to\infty} \operatorname{dist}(q_0(t), \operatorname{co}(q(t))) \leq \frac{1}{k_\alpha N}\sum_{i\in\mathcal{N}} \limsup_{t\to\infty}\|\dot{s}_i(t)\|$;
		\item[\textbf{P2}] $\|q_{ij}(t)\| > d$ for all $t \geq 0$, $i \neq j$;
		\item[\textbf{P3}] The velocity error $\|\dot{q}_i(t) - \dot{q}_0(t)\|$ is ultimately bounded by $\sigma\sqrt{N}/k_\alpha$ where $\sigma = \max_i \sup_t \|\dot{s}_i(t)\|$.
	\end{enumerate}
\end{theorem}

\noindent\textbf{Proof}

\noindent\textbf{Step 1: $s_i\to 0$}
To verify the stability of the closed-loop system, take the derivative of the auxiliary value \( s_i \) \eqref {eq:sliding_var}, we get:
\begin{equation}\label{eq:s_dot}
	\dot{s}_i = \ddot{q}_{i0}+k_\alpha \dot{q}_{i0}+ k_\alpha  q_{i0}-\phi_i.
\end{equation}

Considering the EL system equation \eqref{eq hetero agent} and control law \eqref{eq:control_input}, we have:
\begin{equation*}\label{eq:EL_control}
	M_i(q_i)\ddot{q}_i+C_i(q_i,\dot{q}_i)\dot{q}_i+g_i(q_i)= -k s_i + Y_i(q_{i},\dot{q}_{i},\dot{\zeta}_i,\zeta_i) \hat{\theta}_i. 
\end{equation*}
Combining \eqref{eq:control_input} and \eqref{eq:regression}, we get:
\begin{equation}\label{eq:s_dynamics}
	M_i(q_i)\dot{s}_i+C_i(q_i,\dot{q}_i)s_i= -k s_i + Y_i(q_{i},\dot{q}_{i},\dot{\zeta}_i,\zeta_i) \tilde{\theta}_i, 
\end{equation}
where \(\tilde{\theta}=\hat{\theta}-\theta\).

Choose the Lyapunov function:
\begin{equation*}\label{eq:V1}
	V_1 = \frac{1}{2} s_i^T M_i s_i + \frac{1}{2} (\theta_i - \hat{\theta}_i)^T \Lambda (\theta_i - \hat{\theta}_i).
\end{equation*}
Taking the time derivative of \( V_1 \), and substituting the adaptation law \eqref{eq:adaptation} and equation \eqref{eq:s_dynamics}, we simplify to:
\begin{equation*}\label{eq:V1_dot}
	\dot{V_1} = -k s_i^T s_i \leq 0. 
\end{equation*}

According to Theorem 2.5 in \cite{chen2015stabilization}, \( \dot{V_1} \leq 0 \) indicates that the system is globally asymptotically stable, i.e., the auxiliary value \( s_i \to 0 \).

Moreover, since $\dot{V}_1 = -k\|s_i\|^2$, we have $s_i \in \mathcal{L}_2 \cap \mathcal{L}_\infty$. From the closed-loop dynamics, $\dot{s}_i$ can be shown to be uniformly bounded, which by Barbalat's lemma implies $s_i \to 0$ as $t\to\infty$.



\noindent\textbf{Step 2: Verify the Fencing Property}

From equation \eqref{eq:s_dot}, we have:
\begin{equation}\label{eq:q0_ddot}
	\ddot{q}_{i0}=-k_\alpha \dot{q}_{i0} -k_\alpha {q}_{i0} + \phi_i + \dot{s}_i.
\end{equation}

Choose the Lyapunov function:
\begin{equation*}\label{eq:V2}
	V_2 = \frac{1}{2} \sum_{i\in\mathcal{N}}\sum_{j\in\mathcal{N}}
	\int_{\left\| q_{ij} \right\|}^{\mu}
	\alpha(s)\,ds
	+ \frac{1}{2}k_{\alpha}\sum_{i\in\mathcal{N}}
	\left\| q_{i0} \right\|^2
	+ \frac{1}{2}\sum_{i\in\mathcal{N}}
	\left\| \dot{q}_{i0} \right\|^{2},
\end{equation*}
where $\alpha:(d,\infty) \rightarrow [0,\infty]$ is a suitable function satisfying
\begin{enumerate}
	\item[B1.] When two agents' distance is close to the safety distance, its value tends to infinity  in the sense that  $\lim_{s \to d^+} \alpha(s) = +\infty$;
	\item[B2.] When two agents' distance is not smaller than the communication distance, its value is always zero in the sense that $\alpha_r(s)=0$, $\forall ~ s \in [\mu,+\infty).$
\end{enumerate}

Denote
\begin{equation*}\label{eq:A}
	A=\frac{1}{2} \sum_{i\in\mathcal{N}}\sum_{j\in\mathcal{N}_i}
	\int_{\left\| q_{ij} \right\|}^{\mu}
	\alpha(s)\,ds,
\end{equation*}
\begin{equation}\label{eq:A_dot}
	\dot{A} = \sum_{i\in\mathcal{N}} \frac{\partial A}{\partial q_i}^T \dot{q}_i
    	= -\sum_{i\in\mathcal{N}} \sum_{j\in\mathcal{N}_i} \alpha(\|q_{ij}\|)\frac{q_{ij}^T}{\|q_{ij}\|}\dot{q}_i.
\end{equation}
Since $\sum_{i\in\mathcal{N}}\phi_i = 0$, we have $\dot{A} = -\sum_{i\in\mathcal{N}} \phi_i^T \dot{q}_{i0}$, where $\phi_i$ is defined below.

Taking the time derivative of \( V_2 \), we simplify to:
\begin{equation*}\label{eq:V2_dot}
	\dot{V_2} = \dot{A}+k_{\alpha}\sum_{i\in\mathcal{N}}
	q_{i0}^T\dot{q}_{i0}+\sum_{i\in\mathcal{N}}
	\dot{q}_{i0}^T\ddot{q}_{i0}.  
\end{equation*}
Substitute \eqref{eq:q0_ddot} and \eqref{eq:A_dot},

and
\begin{equation*}\label{eq:phi_i}
	\phi_i= \sum_{j\in\mathcal{N}_i}{\alpha}({\left\| q_{ij} \right\|})\frac{q_{ij}}{\left\| q_{ij} \right\|},
\end{equation*}
then we simplify to
\begin{equation*}\label{eq:V2_dot_final}
	\dot{V_2} =\sum_{i\in\mathcal{N}}\dot{q}_{i0}^T\bigl(-k_\alpha \dot{q}_{i0} + \dot{s}_i\bigr)
	=-k_\alpha\sum_{i\in\mathcal{N}}\left\| \dot{q}_{i0} \right\|^2+\sum_{i\in\mathcal{N}}\dot{q}_{i0}^T\dot{s}_i.
\end{equation*}

Let $\sigma = \max_{i\in\mathcal{N}} \sup_{t\geq 0} \|\dot{s}_i(t)\|$, which is finite since $\dot{s}_i$ is uniformly bounded.
Then
\begin{equation*}\label{eq:V2_bound}
	\dot{V_2} \leq -k_\alpha\sum_{i\in\mathcal{N}}\left\| \dot{q}_{i0} \right\|^2 + \sigma\sum_{i\in\mathcal{N}}\left\| \dot{q}_{i0} \right\|.
\end{equation*}
By Young's inequality $\sigma\|\dot{q}_{i0}\| \leq \frac{k_\alpha}{2}\|\dot{q}_{i0}\|^2 + \frac{\sigma^2}{2k_\alpha}$, we obtain
\begin{equation*}\label{eq:V2_practical}
	\dot{V_2} \leq -\frac{k_\alpha}{2}\sum_{i\in\mathcal{N}}\left\| \dot{q}_{i0} \right\|^2 + \frac{N\sigma^2}{2k_\alpha}.
\end{equation*}
Hence $\dot{V_2} \leq 0$ whenever $\sum_i\|\dot{q}_{i0}\|^2 > N\sigma^2/k_\alpha^2$, which implies that $\|\dot{q}_{i0}\|$ is ultimately bounded by $\sigma\sqrt{N}/k_\alpha$. That is, the velocity error $\|\dot{q}_i - \dot{q}_0\|$ is ultimately bounded, establishing property (P3).



% Proof of (P1) - convex hull fencing
\subsection*{Proof of (P1)}

Define
\begin{equation*}\label{eq:q_tilde}
	\tilde{q}=\frac{1}{N}\sum_{i\in\mathcal{N}}q_{i0},
\end{equation*}
\begin{equation*}\label{eq:q_tilde_dot}
	\dot{\tilde{q}}=\frac{1}{N}\sum_{i\in\mathcal{N}}\dot{q}_{i0},
\end{equation*}
\begin{equation*}\label{eq:q_tilde_ddot}
	\ddot{\tilde{q}}=\frac{1}{N}\sum_{i\in\mathcal{N}}\ddot{q}_{i0},
\end{equation*}

\begin{equation*}\label{eq:q_i0_ddot}
	\ddot{q}_{i0}=-k_\alpha \dot{q}_{i0} -k_\alpha {q}_{i0} + \phi_i + \dot{s}_i.
\end{equation*}

Sum the above expression and then take the average, we get
\begin{equation*}\label{eq:q_tilde_ddot_final}
	\ddot{\tilde{q}}=-k_\alpha \dot{\tilde{q}} -k_\alpha\tilde{q}+\frac{1}{N}\sum_{i\in\mathcal{N}}\dot{s}_i+\frac{1}{N}\sum_{i\in\mathcal{N}}\phi_{i}.
\end{equation*}

As $j\in\mathcal{N}_{i}\Leftrightarrow i\in\mathcal{N}_{j}$ and $q_{ij}=-q_{ji}$, we have
\begin{equation*}\label{eq:sum_phi}
	\sum_{i\in\mathcal{N}}\phi_{i}=0.
\end{equation*}
Then
\begin{equation*}\label{eq:q_tilde_ddot_simplified}
	\ddot{\tilde{q}}=-k_\alpha \dot{\tilde{q}} -k_\alpha\tilde{q}+\frac{1}{N}\sum_{i\in\mathcal{N}}\dot{s}_i.
\end{equation*}

This equation is input-state stable (ISS) as $k_{\alpha}>0$ with input $d(t)=\frac{1}{N}\sum_i\dot{s}_i(t)$. Since $\dot{s}_i$ is bounded, $d(t)$ is bounded. The ISS property gives
\begin{equation*}\label{eq:q_tilde_bound}
	\limsup_{t\to\infty} \|\tilde{q}(t)\| \leq \frac{1}{k_\alpha N}\sum_{i\in\mathcal{N}} \limsup_{t\to\infty} \|\dot{s}_i(t)\|.
\end{equation*}
It, together with the facts that
\begin{equation*}\label{eq:q_tilde_meaning}
	\tilde{q}=\frac{1}{N}\sum_{i\in\mathcal{N}}q_{i0}=\frac{1}{N}\sum_{i\in\mathcal{N}}(q_{i}-q_{o}),
\end{equation*}
indicates that
\begin{equation*}
	\limsup_{t\to\infty}\operatorname{dist}\bigl(q_0(t),\operatorname{co}(q(t))\bigr) \leq \frac{1}{k_\alpha N}\sum_{i\in\mathcal{N}} \limsup_{t\to\infty} \|\dot{s}_i(t)\|,
\end{equation*}
where $\operatorname{dist}(q_0,\operatorname{co}(q)) = \min_{\lambda_i\geq0,\sum\lambda_i=1} \|\sum_{i\in\mathcal{N}} \lambda_i q_i - q_0\|$ denotes the distance from the target to the convex hull of the agents.
Thus, (P1) is proved.

\hfill$\square$

Finally, contradiction is used to prove (P2). If the collision between agents $i$ and $j$ occurs, then $\exists T_0 > 0$ such that $\|q_i(T_0) - q_j(T_0)\| = d$ at the first time for $i \neq j \in \mathcal{N}$. As such

\[
\begin{aligned}
	\int_{\|q_{ij}(t)\|}^{\mu} \alpha(s)\, \mathrm{d}s &= \int_{d}^{\mu} \left(\frac{1}{s - d} - \frac{1}{\mu - d}\right)\, \mathrm{d}s \\
	&= \ln(s - d)\Big|_{d}^{\mu} - 1 = \infty,
\end{aligned}
\]

which implies $V_2(T_0) = \infty$, contradicting the fact that $V_2(t)$ grows at most linearly on any finite time interval (from \eqref{eq:V2_practical}). Thus, the collision between any pair of agents is avoided and (P2) is proved. This completes the proof.